\section{The physical boundary experiment}\label{sec:boundary}
The algebraic extension through $v=0$ in Section~\ref{sec:native} is useful for differentiation, but it does not remove the physical restriction $v\ge0$. We now identify the actual local experiment. The argument is the finite-cell likelihood expansion underlying constrained Gaussian limits \cite{Geyer,vdV}; all normalizations and the synchronized profile are specified here.

Let $u\in\R^{k+1}$ denote the nuisance log-odds and root-sum coordinates. Observe independent categorical samples with $n_j/N\to\lambda_j>0$, and let $\Delta_N=(\Delta_{N,u},\Delta_{N,v})$ be the score at $(u,v)=(0,0)$ divided by $\sqrt N$. Partition
\[
 I=V^{\mathsf T}\diag(\lambda_j\kappa_j)V
   =\begin{pmatrix}I_{uu}&I_{uv}\\I_{vu}&I_{vv}\end{pmatrix},
 \qquad Q=I_{vv}-I_{vu}I_{uu}^{-1}I_{uv}=R^{-1}.
\]
\begin{theorem}[Cone likelihood and synchronized profile]\label{thm:boundary}
For bounded $a\in\R^{k+1}$ and $h\in\R^k_{\ge0}$, the physical local alternatives $u=a/\sqrt N$, $v=h/\sqrt N$ satisfy, uniformly on compact index sets,
\begin{equation}\label{eq:LAN}
 \log\frac{dP_{N,a,h}}{dP_{N,0,0}}
  =(a,h)^{\mathsf T}\Delta_N-\tfrac12(a,h)^{\mathsf T}I(a,h)+o_P(1).
\end{equation}
Under the corresponding alternatives,
\begin{equation}\label{eq:efficient-Y}
 Y_N=R\bigl(\Delta_{N,v}-I_{vu}I_{uu}^{-1}\Delta_{N,u}\bigr)
          \ \Longrightarrow\ N(h,R).
\end{equation}
After quotienting out unrestricted nuisance translation, this is a Gaussian shift with parameter restricted to the physical cone. For synchronized alternatives $h=q_A(\zeta)$, the profiled limiting log likelihood, relative to its value at $h=0$, is
\begin{equation}\label{eq:profile-limit}
 \frac12Y^{\mathsf T}R^{-1}Y-\frac12C_R(Y),\qquad
 C_R(Y)=\inf_\zeta(Y-q_A(\zeta))^{\mathsf T}R^{-1}(Y-q_A(\zeta)).
\end{equation}
The common physical root displacement has scale $N^{-1/4}\zeta$.
\end{theorem}
\begin{proof}
The finite set of clock probabilities is strictly positive on a two-sided neighbourhood of the baseline. Its derivatives through order three are bounded on a smaller neighbourhood. At $u=a/\sqrt N$, $v=h/\sqrt N$, expand each cell log probability. Normalization makes the mean linear term vanish, the expectation of the Hessian is minus the score covariance, and the summed cubic remainder is $O(N^{-1/2})$ on compact index sets. The centred Hessian fluctuation is $o_P(1)$. The finite-dimensional independent-group central limit theorem gives $\Delta_N\Rightarrow N(0,I)$ at the baseline and proves \eqref{eq:LAN}. Under the local alternatives its mean tends to $I(a,h)$ and its covariance to $I$, either by the same cell calculation or the likelihood change of measure. The efficient score in \eqref{eq:efficient-Y} has mean $Qh$ and covariance $Q$, proving the assertion after multiplication by $R$.

In the limiting score experiment the nuisance block is independent of the efficient block. Its mean is $I_{uu}a+I_{uv}h$, which is an unrestricted translation as $a$ varies. Discarding this nuisance translation yields \eqref{eq:efficient-Y}; this is a quotient, not a claim that an unknown nuisance parameter can be simulated without information. Profiling the efficient Gaussian likelihood over $h=q_A(\zeta)$ and completing the square gives \eqref{eq:profile-limit}. The map $q_A$ is proper for full-rank $A$, so the infimum is attained and its optimizing $\zeta$ are bounded whenever $Y$ is bounded. Uniform LAN therefore also gives the iterated compact-localization limit of profiled likelihoods: first restrict $(a,\zeta)$ to a fixed compact set, pass to the limit, and then exhaust the limiting experiment by such sets. No assertion about remote maximizers of an unrestricted finite-sample likelihood is required. Finally $q_A(N^{-1/4}\zeta)=N^{-1/2}q_A(\zeta)$, proving the displacement scale.
\end{proof}
The unrestricted vector $Y$ is a statistic in the limiting experiment, not the constrained maximum likelihood estimator. The latter is a metric projection onto $\R^k_{\ge0}$, or onto $q_A(\R^d)$ in the synchronized experiment. At the boundary it need not be normal; for the nonconvex synchronized image the projection can be set-valued. The distance value in \eqref{eq:profile-limit} is nevertheless unambiguous.

\begin{corollary}[Sharp root-amplitude scale on a physical ray]\label{cor:root-rate}
Fix $a_0\in\R^d\setminus\{0\}$ and $v_0=q_A(a_0)\ne0$. On the one-sided submodel with physical displacement $t a_0$, $0\le t\le T N^{-1/4}$, and nuisance $u=a/\sqrt N$ in a fixed bounded set, the minimax squared risk for estimating $t$ is of order $N^{-1/2}$, for each fixed $T>0$. Thus the amplitude scale is $N^{-1/4}$, whereas its squared amplitude is on the $N^{-1/2}$ parameter scale.
\end{corollary}
\begin{proof}
The variance parameter is $v=t^2v_0$. Set $J_0=v_0^{\mathsf T}Qv_0>0$ and use the baseline efficient score to define
\[
 \widehat b_N=\frac{v_0^{\mathsf T}QY_N}{\sqrt N J_0},\qquad
 \widehat t_N=\sqrt{\max\{\widehat b_N,0\}}.
\]
The finite-cell moment expansion used above gives, uniformly in the stated local set,
$\mathbb E(\widehat b_N-t^2)^2\le C/N$; the mean error is $o(N^{-1/2})$ and the variance is $O(N^{-1})$. For $b\ge0$,
$(\sqrt{\max\{x,0\}}-\sqrt b)^2\le|x-b|$.
Cauchy--Schwarz therefore gives the upper bound $C'/\sqrt N$.

For the lower bound take the two physical alternatives $t=0$ and $t=\tau N^{-1/4}$ with the same nuisance, where $0<\tau\le T$ is sufficiently small. Their likelihood divergence is bounded by a constant times $\tau^4$, by \eqref{eq:LAN} or its finite-cell Taylor proof. Their total variation is consequently bounded away from one. The two-point testing bound for squared loss gives a positive constant times the squared separation $\tau^2N^{-1/2}$. The claim concerns this labelled nonnegative ray; it does not assert identifiability of the sign of a general loading vector.
\end{proof}

\section{Direct count-to-contact experiment comparison}\label{sec:score-contact}
Here the observations differ from fixed-count categorical marks. Let the mark laws, hence $\kappa_j$, be known, and observe independent exposures
\[
 N_j\sim\operatorname{Pois}(N\lambda_j),\qquad j=1,\ldots,m.
\]
This is the known-mark-law subexperiment of an observed-exposure protocol. Fix an interior centre $\lambda^0>0$. Write $D$ for the derivative of the native moment functional at that centre:
\[
 D_{\cdot j}=-\frac{\rho_j}{(\lambda_j^0)^2\kappa_j}w_{2n}(T_j),
 \quad \Lambda=\diag(\lambda_j^0),\quad
 \Sigma=D\Lambda D^{\mathsf T}>0.
\]
The positivity follows from Vandermonde rank. For a local moment index $h\in\R^p$ in a fixed compact ball, define the least-favourable target submodel
\begin{equation}\label{eq:least-favourable}
 \lambda_N(h)=\lambda^0+N^{-1/2}B_*h,
 \qquad B_*=\Lambda D^{\mathsf T}\Sigma^{-1}.
\end{equation}
It has $\theta(\lambda_N(h))=\theta(\lambda^0)+h/\sqrt N+O(N^{-1})$. Its local information is $\Sigma^{-1}$, since $DB_*=I_p$ and $B_*^{\mathsf T}\Lambda^{-1}B_*=\Sigma^{-1}$.

Let $\LL$ be the known normal-compression operator, possibly rank deficient, and put $r=\rank\LL$. Choose an orthonormal coordinate frame $F$ for its range. Generate independent $U_j\sim\operatorname{Unif}(-1/2,1/2)$, independently of the counts, and retain only
\begin{equation}\label{eq:direct-statistic}
 T_N=F^{\mathsf T}\LL D\frac{N_\cdot-N\lambda^0+U}{\sqrt N}.
\end{equation}
The randomization law is part of the protocol; its realized auxiliary values are not retained. This statistic is formed directly by fixed linear weights on the counts. No estimator of $\theta$ is constructed first. The small randomization is essential for an unambiguous statistical reduction: an exact real-valued linear encoding of an integer vector can be injective even when its real matrix has deficient rank. Weak convergence alone would not exclude that lattice encoding.

\begin{lemma}[Jittered Poisson local comparison]\label{lem:poisson-TV}
Uniformly for $h$ in a fixed compact set, the distribution of
\[
 X_N^\circ=(N_\cdot-N\lambda^0+U)/\sqrt N
\]
under \eqref{eq:least-favourable} tends in total variation to
$N(B_*h,\Lambda)$. The full jittered vector and the raw counts are exactly equivalent experiments, by independent jittering and rounding. The raw local target experiment is therefore asymptotically equivalent to $Z\sim N(h,\Sigma)$.
\end{lemma}
\begin{proof}
For one count with mean $N\lambda$, jittering gives a step density with cell width $N^{-1/2}$. Stirling's formula, uniformly for $\lambda$ in a fixed positive compact interval, shows that on $|x-\sqrt N(\lambda-\lambda^0)|\le N^{1/12}$ its ratio to the normal density of mean $\sqrt N(\lambda-\lambda^0)$ and variance $\lambda$ is
\[
 1+O\bigl(N^{-1/2}(1+N^{1/4})\bigr)=1+o(1).
\]
The variation of the normal density inside a jitter cell has the same vanishing relative bound. Outside this region the probabilities tend to zero uniformly, for example by the second-moment bound. Thus the density difference tends to zero in $L^1$. Apply this in the fixed number of independent coordinates. Since $\lambda_N(h)-\lambda^0=O(N^{-1/2})$ uniformly, replacing the diagonal normal covariance by $\Lambda$ changes total variation by $o(1)$.

Rounding $\sqrt N X_N^\circ+N\lambda^0$ to the nearest integer recovers each count almost surely, while independent jittering gives the reverse Markov kernel. In the limiting normal vector $X\sim N(B_*h,\Lambda)$, the statistic $Z=DX$ has law $N(h,\Sigma)$ and is sufficient: conditional simulation is
\[
 X=B_*Z+E,\qquad
 E\sim N(0,\Lambda-\Lambda D^{\mathsf T}\Sigma^{-1}D\Lambda),
\]
with $E$ independent and its law independent of $h$. These kernels, the total variation approximation, and contraction of total variation under kernels prove the experiment comparison uniformly on the compact parameter set. In particular this proves the asserted local Le Cam comparison, not merely a central limit theorem.
\end{proof}

\begin{theorem}[Exact local information retained by contact scores]\label{thm:contact-experiment}
The experiment retaining \eqref{eq:direct-statistic} is, uniformly on compact local parameter balls, asymptotically equivalent to
\begin{equation}\label{eq:contact-limit}
 T\sim N(F^{\mathsf T}\LL h,\ F^{\mathsf T}\LL\Sigma\LL^{\mathsf T}F).
\end{equation}
Its limiting information matrix for $h$ is
\begin{equation}\label{eq:contact-information}
 I_{\rm contact}=\LL^{\mathsf T}(\LL\Sigma\LL^{\mathsf T})^\dagger\LL
 =\Sigma^{-1/2}P_{\ran(\Sigma^{1/2}\LL^{\mathsf T})}\Sigma^{-1/2}.
\end{equation}
It has rank $r$ and kernel $\ker\LL$. The contact-score reduction is locally asymptotically equivalent to the raw target subexperiment if and only if $\LL$ is injective. Hence it retains all local moment information exactly when the within-contact polynomial products span $V_{2n}$.
\end{theorem}
\begin{proof}
Apply the fixed linear map $F^{\mathsf T}\LL D$ to the total variation approximation in Lemma~\ref{lem:poisson-TV}. Contraction gives convergence in total variation to \eqref{eq:contact-limit}, whose covariance is positive definite in its $r$ coordinates. Identity kernels therefore compare these experiments in both directions with error tending to zero. The usual normal-shift likelihood gives the first expression in \eqref{eq:contact-information}. For $M=\LL\Sigma^{1/2}$, the identity $M^{\mathsf T}(MM^{\mathsf T})^\dagger M=P_{\ran M^{\mathsf T}}$ gives the second. Its rank and kernel follow immediately.

If $\LL$ is injective, $Z$ is recovered from $\LL Z$ by a left inverse. Thus \eqref{eq:contact-limit} is equivalent to $N(h,\Sigma)$, and Lemma~\ref{lem:poisson-TV} gives equivalence to the raw target subexperiment. Conversely choose distinct $h_0,h_1$ in a local parameter ball with $h_1-h_0\in\ker\LL$. Their limiting contact distributions coincide, whereas $N(h_0,\Sigma)$ and $N(h_1,\Sigma)$ have strictly positive total variation distance. By the triangle inequality, any kernel purporting to reconstruct both raw limit laws has error at least half that distance. The same positive lower obstruction holds asymptotically by the uniform comparisons. The polynomial criterion is Theorem~\ref{thm:fibre}.
\end{proof}
The uncentred contact block is $B_\nu(\theta)=\LL_\nu\theta$. Its local fluctuation is precisely the linear type appearing in \eqref{eq:direct-statistic}. Passing to inverse-Hessian contact scores applies the fixed invertible block transformation $E_\nu\mapsto-B_{\nu,0}^{-1}E_\nu B_{\nu,0}^{-1}$, so it preserves the information statement. These are correlated raw-data score statistics, not independently noisy observations of distance values.

The centre $\lambda^0$ and the local submodel \eqref{eq:least-favourable} are specified before observing the data. The theorem does not claim global sufficiency when both exposure and mark laws are unknown, or that a plug-in estimate of the centre automatically gives the same comparison. It supplies a genuine direct reduction and a complete local information-loss calculation in the stated full-dimensional target experiment. When $\LL$ is injective, exact range-compatible score encoding does not introduce a $\gamma^{-2}$ statistical variance penalty; independent measurement errors and numerical perturbations have the distinct stability behaviour of Theorem~\ref{thm:stability} and Appendix~\ref{app:offset}.
